This section describes the implementation of \GMR and \SOE models in
\GAMS using the \GTAP12 dataset. We first describe the content of the
distribution directory, and provide the logic of the model
buildstream. We then discuss a number of issues related to data
manipulation, namely aggregation across commodities and regions, data
filtering and re-balancing. Finally, we discuss calibration of
elasticities for international trade in \GMR and \SOE models,
providing an intuitive (Marshallian) argument on how models can be
parametrized in order to approximate the same responsiveness to
policy shocks.

\subsection{Distribution Folders}

In this section we overview the structure of the distribution
directory. Practicalities on how to run the code is provided in
Appendix \ref{installation}. Details of the directories and \GAMS
programs provided in \GTAPINGAMS version 12 are provided in Appendix
\ref{programs}. More detailed comments may also be found interspersed
in the GAMS code presented in Appendix \ref{code}.

The distribution directory contains seven second-level subdirectories,
whose content we now overview in alphabetical sequence:

\begin{itemize}

\item[\textsc{build}] contains a command script which automates
  dataset construction and testing. This script (\texttt{build.gms})
  calls programs stored in the \textsc{code}
  subfolder. \textsc{filter.gms}, a program which filters and then
  rebalances the core dataset (see data filtering section below),
  \textsc{gtap12data.gms}, a program which reads a \textsc{.gdx}
  \GTAPINGAMS dataset, \textsc{gtapaggr.gms}, a program which performs
  dataset aggregation, and \textsc{cdecalib.gms}, a program which
  computes coefficients of the \CDE demand system which are consistent
  with price and income elasticities of demand.

The distributed copy of \texttt{build.gms} reads, filters and
aggregates data files for 2011. At the end of the computational
process, listing files can be found in the
\textsc{build} directory, while the associated \textsc{.gdx} data files
generated by these routines are saved in the \textsc{gamsdata}
directory.

This directory also contains programs for consistency testing of
generated datasets. \texttt{test.gms} runs \texttt{bmkchk.gms},
\texttt{cdechk.gms} and \texttt{mcpmge.gms}. See Appendix
\ref{programs} for a description of these programs.

\item[\textsc{code}] contains the main \GAMS program files which are
included in this distribution. These files need not reside in this
directory, but they are self contained and should always be moved as a
collection. Two of the files in this directory contain \GAMS code for
the canonical static model, both specified as a mixed complementarity
problem in two alternative representations:

\begin{description}
\item{\textsc{mge.gms}} The standard model with tabular \GAMS/\MPSGE
representation, and
\item{\textsc{mcp.gms}} The standard model with an algebraic \GAMS/\MCP
representation.
\end{description}

\item[\textsc{defines}] contains mapping files for aggregations of
\GTAP 12 datasets. Files ending with \textsc{.map} define an
aggregation in terms of the source dataset and mappings from sets in
the source to sets in the target. Three single dimensional sets and
three two-dimensional tuples are included in each mapping file. The
sets and tuples are defined as follows:
\begin{Verbatim}
set rr          Regions in the aggregation,
    ii          Commodities in the aggregation,
    ff          Primary factors in the aggregation,
    mapr(r,rr)  Mapping of source to target regions (from-to),
    mapi(i,ii)  Mapping from source to target commodities
    mapf(f,ff)  Mapping from source to target factors;
\end{Verbatim}

As an example, here is a set and mapping which aggregates skilled
labor types:
\begin{Verbatim}
set ff  Factors in the aggregated data /
    skl     Skilled workers,
    lab     Agricultural and unskilled workers,
    lnd     Land,
    cap     Capital,
    res     Natural resources /;

set mapf(f,ff)  Factor aggregation assignments /
    mgr.skl Officials and mangers (ISCO-88 Major Groups 1-2),
    tec.skl Technicians and associate professionals
    clk.skl Clerks
    srv.lab Service and market sales workers
    lab.lab Agricultural and unskilled workers (Major Groups 6-9)
    lnd.lnd Land,
    cap.cap Capital,
    res.res Natural resources /;
\end{Verbatim}

Four alternative regional / commodity aggregation are included in the
\textsc{defines} directory to illustrate how aggregation works in
\GTAPINGAMS. The first is the full disaggregate dataset
\textsc{gtapingams.map} which retains all regions and markets from the
\GTAP 12 Data Base after original data have been filtered and a few sector
labels have been changed.\footnote{The sectoral identifier which
differ from the \GEMPACK model are as follows: \textsc{oil} (crude
oil) is relabeled \textsc{cru}, \textsc{coa} (coal) becomes
\textsc{col}, \textsc{p\_c} (petroleum and coal products) becomes
\textsc{oil}, \textsc{ely} (electricity) becomes \textsc{ele}, and
\textsc{ele} (electronic equipment) becomes \textsc{eeq}.}

Three aggregated datasets are based on regional disaggregation which
distinguishes the G20, other oil exporters, a composite region
representing other low income countries and a composite region
representing other middle income countries. The \textsc{g20\_iea}
aggregates commodities into 23 sectors included in the International
Energy Agency energy database, and \textsc{g20\_macro} aggregates to a four
commodity sectoral mapping distinguishing agriculture, manufacturing,
services and energy.

\item[\textsc{doc}] contains the paper you are reading (please check
the version).

\item[\textsc{forensics}] contains \GAMS programs illustrating how the
distribution can be used to carry out policy analysis. Specifically,
we provide code to evaluate the sensitivity of unilateral trade tax
reform with respect to final demand system and model closure. Program
\textsc{ssagen.gms} generates \textsc{ssasolve.gms} and
\textsc{ssamerge.gms}, \GAMS scripts which
invokes a set of jobs in parallel and then merges the model results
into a pivot data format. The jobs provide a sequence of
calculations under alternative closures and policy region, each
corresponding to an invocation of the scenario file
\textsc{tgrid.gms}. Results from the sequence of simulations are
merged into a PivotChart dataset in an Excel workbook.

\item[\textsc{gamsdata}] contains the datasets generated by running
\texttt{build.gms} in the \textsc{build} directory. More specifically,
files that are created and stored in this directory have a GAMS
data-exchange (\textsc{*.gdx}) extension and are constructed data
files that are either used as the source data for aggregation (e.g.
\textsc{gsd.gdx}) or the target data emerging from the data
aggregation (e.g. the filtered \GTAP data stored as
\textsc{gtap12ingams.gdx}), data files which are then then used in
conjunction with the model files). Here are the sizes of the 2011
datasets produced for a 2011 base year:
\begin{Verbatim}
  gsd.gdx                 43,910,349    ! From flex2gdx.gms
  gsd_5.gdx               26,690,775    ! From filter.gms -- nd=5
  gtap12ingams.gdx         25,923,318    ! From gsd_5
  g20.gdx                  3,770,394    ! From gtap12ingams
  g20_iea.gdx                992,710    ! From g20 (23 sectors)
  g20_macro.gdx              113,224    ! From g20 (4 sectors)
\end{Verbatim}

\item[\textsc{gtapdata}] contains \texttt{.zip} files from the \GTAP
distribution (GDX package). These files are not provided in the
\GTAPINGAMS distribution, but must be copied here by the user. They
are shipped in a file that includes all reference years:
\begin{Verbatim}
  GDX12a_AY16994.zip
\end{Verbatim}

\end{itemize}

\subsubsection{Aggregation across commodities}

Any \GTAPINGAMS dataset may be aggregated into fewer regions, sectors
and primary factors. This permits a modeler to do preliminary model
development using a small dataset to ensure rapid response and a short
debug cycle. After having implemented a small model, it is then a
simple matter to expand the number of sectors and/or regions in order
to obtain a more precise empirical estimate.

Conveniently, all \GTAP datasets are defined in terms of three primary
sets: $i$, the set of sectors and produced commodities, $r$ the set of
countries and regions, and $f$ the set of primary factors.
Table~\ref{allsectors} presents the identifiers for the 65 \GTAP 12
sectors (plus CGD investment goods) in their most disaggregate form
(see Appendix \ref{sectors}). These sectors may be aggregated freely to produce
more compact dataset.

Aggregation across commodities involves the reclassification of goods
at \textit{a benchmark point in the price space}. If we index the
\textit{disaggregate} goods using $i,j$ and index the corresponding
\textit{aggregated} goods with $I,J$, we can identify commodities
combined to form $I$ as $i \in \mathcal{I}_I$. Aggregation at
reference prices $\bar p_i$ implies that the value of aggregate demand
equals the value of the constituent disaggregate
demand:\footnote{\label{harberger}We use $\bar p_i$ and $\bar p_I$ in
this expression to emphasize the role of price as an aggregation
weight, but without loss of generality, we may assume that the
reference price vectors $\bar p_i = 1 \quad \forall i$ and $\bar p_I =
1 \quad \forall I$. This Harberger normalization greatly simplifies
notation. When units are scaled so that the unit cost at this point
is one, values shares at this point are identically equal to share
parameters in the \CES function, i.e. $\alpha_i = \theta_i$, and
benchmark demands are equal to value shares, i.e. $\bar C_i =
\theta_i.$}
\begin{equation}
\label{aggregation} \bar p_I C_I(\bar p) = \sum_{i \in \mathcal{I}_I}
\bar p_i C_i(\bar p). \end{equation} Furthermore, \textit{consistent
exact aggregation} implies that both demand and the elasticity of
demand at the benchmark point in the aggregate model coincides with
the demand response in the disaggregate model, taking benchmark prices
into account. Formally:
\begin{equation}
\label{ownpriceconsistency}
\bar p_J \left. \frac{\partial C_I}{\partial p_J} \right|_{\bar p} = \sum_{i \in \mathcal{I}_I, j \in
\mathcal{I}_J} \bar p_j \left. \frac{\partial C_i}{\partial p_j}
\right|_{\bar p}.
\end{equation}
An implication is that the Allen-Uzawa elasticity of substitution
(\AUES) matrix of the aggregate model has to correspond to
value share-weighted average of terms in the \AUES matrix of the
disaggregate model:
\begin{equation}
\label{exactform}
\hat \sigma_{IJ} = \sum_{i\in \mathcal{I}_I, j \in \mathcal{I}_J}
\frac{\theta_i}{\hat \theta_I} \frac{\theta_j}{\hat \theta_J}
\sigma_{ij}
\end{equation}
where $\hat \theta_I, \hat \theta_J$ are the value share of inputs $I$
and $J$ in the aggregate cost function, and $\hat \sigma$ is the \AUES
matrix in terms of the aggregate prices.\footnote{Note that exact
aggregation based on (\ref{exactform}) is only possible when working
with a specific functional form, in which case own-price elasticities
in the aggregated model are:
\begin{equation}
\label{exactagg}
 \epsilon_{I} = \hat \theta_I \sigma_{II}  = \hat \theta_I \sum_{(ij) \in \mathcal{I}_I}
\frac{\theta_i}{\hat \theta_I} \frac{\theta_j}{\hat \theta_I} \sigma_{ij}
\end{equation}
Note also that if we had information about benchmark value shares
($\theta_i$) and benchmark elasticities of demand ($\epsilon_i$), but
no information regarding off-diagonal terms in the \AUES matrix
$\sigma_{ij},\, i\neq j$, the vectors $\theta_i$ and $\epsilon_i$ may
or may not be consistent with a given cost function. In this setting,
aggregation can use a simplifying assumption that the average
off-diagonal \AUES in the aggregation-relevant terms equals the
overall average off-diagonal value for each of the constituent goods.
From this it follows that average share-weight off-diagonal AUES value
associated with input $i$ is given by:
$$ \bar \sigma_{i} = \frac{\epsilon_i}{\theta_i - 1}$$
Substituting $\bar \sigma_i$ for $\sigma_{ij}$ in (\ref{exactagg}),
we have a
formula for an inexact aggregation:
\begin{equation}
\label{theformula}
\hat \epsilon_I = \sum_{i \in \mathcal{I}_I} \left( \frac{1-1/{\hat
\theta_I}}{1 - 1/\theta_i} \right) \frac{\theta_i}{\theta_I}
\epsilon_i
\end{equation}
In either exact aggregation (working from a full \AUES matrix) or
inexact aggregation (based on average off-diagonal values), we find
that $\theta_i < \theta_I \quad \forall i \in \mathcal{I}_i$,
hence the elasticity of demand for composite good $I$ is less
than a weighted sum of the elasticities of demand in the
constituent goods:
$$|\hat \epsilon_I| < \sum_{i \in \mathcal{I}_I} \frac{\theta_i}{\hat \theta_{I}}
|\epsilon_i|\,.$$}

In practice, our aggregation routine proceeds as follows. First, we
construct the full \AUES matrix from the \CDE parameters at the
benchmark point ($\texttt{subpar}_{ir}$ in the \GTAP Data Base), using
expression (\ref{sigmacde}) in Appendix \ref{CDE}. Second, to ensure a
consistent demand response when the data is aggregated across
commodities, we use the aggregated value shares and the diagonal terms
of the \AUES matrix to recalibrate coefficients of the \CDE demand
system.

\subsubsection{Aggregation across regions}

Suppose that we aggregate across regions as well as across
commodities.\footnote{Note that regional identifiers in the full
dataset correspond to ISO alpha-3 country codes. Users can define
their own aggregations of the \GTAP Data Base and use any labels to
describe regions. For technical reasons, if a dataset is to be
used with \MPSGE, then \textit{regional identifiers can have at most 4
characters}. Table~\ref{factors} presents the three-character
identifiers which are normally used for primary factors. } Cost
functions in regions $r$ and $r'$ are disjoint, so the demand response
for $p_{ir}$ does not depend on $p_{ir'}$. We therefore begin by
calculating price elasticities for the aggregate goods within each
region, and then produce a quantity-weighted aggregation of
elasticities across regions. Therefore, the own-price elasticity of
demand for aggregate commodity $I$ in aggregate region $R$ is given
by:
$$\epsilon_{IR} = \frac
{\sum_{r\in \mathcal{R}_R} \epsilon_{Ir} \bar C_{Ir}}
{\sum_{r\in \mathcal{R}_R}               \bar C_{Ir}}\,,$$
where $\bar C_{Ir}$ is benchmark consumption. As noted in footnote
\ref{harberger}, because of the Harbgerger normalization, it is also
equal to the value shares.

Own-price elasticities together with value shares define diagonal
terms in the Allen-Uzawa elasticity matrix. To compute off-diagonal
\AUES terms, we use a value-weighted aggregation of cross-price
elasticities across regions. In particular, cross-price (compensated)
elasticities are given by:
$$\epsilon_{IJr} = \frac
{\sum_{i\in \mathcal{I}_I, j \in \mathcal{I}_J} \sigma_{ijr}
\bar C_{ir} \bar C_{jr}}
{\sum_{i\in \mathcal{I}_I, j \in \mathcal{I}_J}  \bar C_{ir} \bar
C_{jr}} ~ \theta_{Jr}. $$
To ensure that these terms are consistent with the adding up condition
of the \AUES matrix, we average these terms across
constituent regions and scale them appropriately:
$$\sigma_{IJR} = \lambda_{IR} \left( \frac{\sum_{r \in \mathcal{R}_R}
\epsilon_{IJr} ~ \psi_{IJrR}} {\theta_{JR}} \right) ~  \quad \forall
I \neq J\,,$$
in which
$$\psi_{IJrR} = \frac{\bar C_{Ir} \bar C_{Jr}}
{\sum_{r'\in \mathcal{R}_R} \bar C_{Ir'} \bar C_{Jr'}}\,,$$ and the
scaling factor $\lambda_{IR}$ is chosen so that $\sum_J \sigma_{IJR} ~
\theta_{JR} = 0$. Given our somewhat ad-hoc treatment of off-diagonal
terms, it seems sensible to define aggregations from the \GTAP Data Base
with full commodity and regional details rather than from a dataset
which has already been aggregated.

\subsection{Data filtering}

The \GTAP source data in original form presents substantial challenges
for calibrated models processed using direct solution methods (e.g.,
\textsc{path}, \textsc{conopt} or \textsc{ipopt}). In our experience,
most numerical problems with \GTAPINGAMS models can be traced to
density of the source data in which we find large numbers of small
coefficients. These coefficients portray economic flows which are a
negligible share of overall economic activity, yet impose a
significant computational burden during matrix factorization.

In order to ``filter'' these economically insignificant value and
reduce dimensionality of the problem, \GTAPINGAMS includes a \GAMS
program (\textsc{filter.gms}) which removes small values which are
smaller than a specified tolerance level. An input to this program
(\textsc{nd}) determines the filter tolerance, i.e. the number of
decimals for the smallest coefficient to be retained in the data. For
example, when $nd=4$, the smallest coefficient in the benchmark social
accounting matrix is 0.0001. When $nd=6$, the smallest number is
0.000001, etc. Larger values of \textsc{nd} retain a larger number of
small coefficients in the filtered dataset. Filtering makes a \GTAP
Data Base smaller, as illustrated in Table \ref{nonzeros} in which it
can be seen that filtering reduces the size of a \GTAP Data Base by
somewhere between 20\% and 50\%, depending on the filtering tolerance.

\begin{table}[!t]
\centering
\caption{Filtering results for entire dataset (2011 data).}\label{nonzeros}
\small
\begin{tabular}{rc|r|rrr|>{\raggedleft}p{0.3in}>{\raggedleft}p{0.3in}>{\raggedleft\arraybackslash}p{0.3in}}
\toprule
          & \textit{} & \multicolumn{4}{c}{Nonzeros (1000)} \\
          & \textit{} &  Dataset & \multicolumn{3}{c}{After Filtering} & \multicolumn{3}{|c}{\% Reduction } \\
\addlinespace
\midrule
\addlinespace
          &  &      $nd=$ & $4$     & $5$     & $6$
          & $4$     & $5$     & $6$ \\
    Domestic Intermediate & \textit{vdfm} & 478.8 & 233.3 & 290.7 & 342.2 & 51    & 39    & 29 \\
    Imported Intermediate & \textit{vifm} & 478.8 & 186.5 & 247.6 & 300.3 & 61    & 48    & 37 \\
    Bilateral trade & \textit{vtwr} & 776.4 & 175.4 & 314.2 & 463.8 & 77    & 60    & 40 \\
    Transport margins & \textit{vxmd} & 109.8 & 400.9 & 591.1 & 766.8 & 63    & 46    & 30 \\
\addlinespace
\bottomrule
\end{tabular}%
\vspace{0.4in}
\end{table}%

Most of the reduction in non-zeros results from the elimination of
small intermediate inputs and bilateral trade flows. Importantly, the
filtering procedure has differential impacts on regions and markets in
the database, and the choice of a filter tolerance depends to some
extent on the size of the countries to be included in a given
analysis. For our analysis of the G20 countries in the present paper,
a filter tolerance $nd=5$ seems suitable, whereas $nd=7$ leads to
numerical problems due to small coefficients. Therefore, depending on
the nature of the policy question under investigation, different
filtering strategies may be adopted. Individual researchers may have
their own opinions about how to select parameter values should be
rounded to zero. The version of \texttt{filter.gms} provided with the
\GTAPINGAMS distribution is intended to provide a starting point for
this step in the dataset development process.

Removing small entries from the dataset implies that the resulting
filtered dataset no longer represent a micro-consistent matrix.
However, unlike earlier versions of \GTAPINGAMS, we do not use a
nonlinear optimization framework to rebalance the data. Instead,
\texttt{filter.gms} moves imbalances resulting from omitted
coefficients into either factor supplies or investment demand
depending on the sign of imbalance which appears following filtering.
This approach to reconciliation is simple to implement provided that
the inconsistencies resulting from filtering are small. The
reconciliation methodology implemented in \texttt{filter.gms} is
perhaps less useful for large scale recalibration exercises such as
might arise with a wholesale change in benchmark tax rates or production
structure.

In practice, we note that for $nd=5$, the data rebalancing procedure
imposes very small changes in the remaining nonzero elements, but it
substantially reduces the number of coefficients in the dataset as
indicated in Table \ref{nonzeros}. Nearly half of the trade flow and
imported intermediate demand coefficients can be dropped, resulting in
a reduction in the size of the dataset from nearly 44 MB to 27 MB.
More importantly, sparsity tends to improve robustness. Larger models
are easier to solve when the underlying datasets have been carefully
filtered.

\subsection{Calibration of trade elasticities}

Our \SOE model is based on the idea that goods are differentiated both
on the basis of origin and destination. In this setting, production
supplies goods for the domestic and export markets with a technology
capturing the idea that the quantities supplied to these submarkets
depend on relative prices. Each traded good then has three
associated trade elasticities: i) $\texttt{etadx(i)}$, the elasticity
of transformation between goods produced for the the domestic and
export markets; ii) $\texttt{esubdm(i)}$, the elasticity of
substitution between domestic and imported varieties entering
intermediate and final demand, and iii) $\texttt{esubm(i)}$, the
elasticity of substitution between imports from different countries.

In the \GMR model corresponding to the standard \GTAP model,
$\eta^{DX} = \infty$, so goods produced for the domestic and export
markets are perfect substitutes. In the \GTAPINGAMS framework,
$\eta^{DX}$ can take on any nonnegative value.  A finite value is
typically specified in the \SOE model. When doing this, we adjust the
domestic-import demand elasticity to maintain \textit{consistency}
with trade responses in the standard \GTAP model. We interpret
consistency from a Marshallian perspective. In a given market, suppose
the import demand elasticity is $\epsilon$,\footnote{When value shares
are small, the price elasticity of demand for imports is approximately
equal to $-\texttt{esubdm(i)}$.} and benchmark rate of protection is
$\bar t$. Scaling units so that the FOB supply price is $1$, the
demand price is $1+\bar t$, and the quantity demanded in the benchmark
is $1$.

When the import supply elasticity from partner countries, $\eta$,
is infinite, the import supply price is fixed  ($p=1$).
Therefore, if the tariff was removed and demand is isoelastic, the
new equilibrium quantity and price (in an isoelastic model) are given by:
\begin{equation}
\label{newqty}
q^* = \left(1 + \bar t \right)^{|\epsilon|}\,.
\end{equation}
With a fixed import supply price, the quantity demanded at $p=1$ is
the new equilibrium. This outcome is illustrated in Figure
\ref{etainfty}, in which benchmark import demand is at price-quantity
point $a$, benchmark import supply is at point $b$, and there is a
benchmark trade barrier equal to 20\%. The free-trade equilibrium is
at point $c$, reflecting benchmark import demand elasticity $\epsilon = 5$.

\begin{figure}[p!]

\begin{subfigure}[b]{0.9\textwidth}
\centering
\includegraphics[height=0.4\textheight]{figures/figure9a.tikz}
\caption{ $\epsilon=5$, $\eta=\infty$.}
\label{etainfty}
\end{subfigure}

\vspace{0.2in}

\begin{subfigure}[b]{0.9\textwidth}
\centering
\includegraphics[height=0.4\textheight]{figures/figure9b.tikz}
\caption{$\epsilon=16.7$, $\eta=7.1$.}
\label{theta30}
\end{subfigure}

\caption{Consistent Calibration of $\epsilon$ and $\eta$.}

\end{figure}

Now instead of a perfectly elastic import supply we consider $\eta <
\infty$. In this case, removing tariffs implies a new equilibrium
supply price $p^*$, and the associated equilibrium supply and demand
$q^*$ from (\ref{newqty}) is given by:
$$q^* = \left( \frac{1+\bar t}{p^*} \right)^{|\epsilon|} =
\left( p^* \right)^\eta\,.$$
The aggregate trade response under $\eta < \infty$ equals the
response with $\eta = \infty$ when the demand elasticity ($\epsilon$)
is adjusted upward:
$$p^* = (1 + \bar t)^{|\epsilon|/\eta}\,,$$
the reparametrized elasticity of demand is given by:
$$|\epsilon|  = \frac{|\epsilon|}{1 - |\epsilon|/\eta}\,.$$

Let $\theta = 1 - |\epsilon|/\eta$. Consistency of trade quantity
responses in the revised model implies that $\theta \geq 0$. When
$\eta = \infty$, then we have the original model with a horizontal
import supply curve. Setting $\theta=0.3$, the new Marshallian model
appears in Figure \ref{theta30}. As in the previous model, $a$ and $b$
are benchmark demand and supply with a 20\% tariff barrier. While free
trade in the original model is at point $c$, free trade in the revised
model is at point $d$. Supply is less elastic, so social surplus
increases are captured by the supply side of the market, and the
equilibrium price following tariff removal increases by over 10\% from
the benchmark level.

In the \GAMS code, this adjustment takes place at the top of
\texttt{gtap12data.gms} file. The run-time environment variable
\texttt{thetadx} defines how the adjustment of export supply and
import demand is to be calibrated to match the \GMR model. The
associated assignments are reported here.

\begin{Verbatim}
parameter       etadx(g)        Export elasticity of transformation D vs E,
                esubdm(i)       Elasticity of substitution D vs M,
                vxm(g,r)        Export market supply;

$if not set thetadx $set thetadx 0.5

etadx(i)$(1-%thetadx%) = esubd(i)/(1-%thetadx%);
etadx(i)$(not (1-%thetadx%)) = +inf;
esubdm(i) = esubd(i)/%thetadx%;
vxm(i,s) = (vst(i,s) + sum(r,vxmd(i,s,r)))$(1/etadx(i));
\end{Verbatim}

%%% Local Variables:
%%% mode: LaTeX
%%% TeX-master: "documentation"
%%% End:
